\documentclass[12pt,a4paper]{report}

\usepackage[margin=1in]{geometry}
\usepackage[T1]{fontenc}
\usepackage[utf8]{inputenc}
\usepackage{lmodern}
\usepackage{setspace}
\usepackage{microtype}
\usepackage{graphicx}
\usepackage{svg}
\usepackage{booktabs}
\usepackage{longtable}
\usepackage{array}
\usepackage{amsmath}
\usepackage{amssymb}
\usepackage{siunitx}
\usepackage{hyperref}
\usepackage{xcolor}
\usepackage{listings}
\usepackage{caption}
\usepackage{float}
\usepackage{tikz}
\usepackage{enumitem}
\usepackage{fancyhdr}
\usepackage{pdflscape}
\usepackage{placeins}
\usepackage{needspace}
\usepackage{natbib}

\usetikzlibrary{arrows.meta, positioning, shapes.geometric, fit, calc}
\onehalfspacing
\setlength{\parskip}{0.45em}
\setlength{\headheight}{15pt}
\pagestyle{fancy}
\fancyhf{}
\fancyhead[L]{Global MultiModalSoilNet Part 2}
\fancyhead[R]{Technical Report}
\fancyfoot[C]{\thepage}

\hypersetup{
    colorlinks=true,
    linkcolor=blue!60!black,
    urlcolor=blue!60!black,
    citecolor=blue!60!black,
    pdftitle={Global MultiModalSoilNet Report},
    pdfauthor={Mayur Kunnitigi Matada}
}

\lstdefinestyle{pycode}{
    language=Python,
    basicstyle=\ttfamily\scriptsize,
    keywordstyle=\color{blue!70!black}\bfseries,
    commentstyle=\color{green!40!black},
    stringstyle=\color{orange!70!black},
    numbers=left,
    numberstyle=\tiny\color{gray},
    numbersep=6pt,
    frame=single,
    breaklines=true,
    showstringspaces=false
}

\newcommand{\feat}[1]{#1\par}
\newcommand{\codepath}[1]{\texttt{\small \detokenize{#1}}}
\newcommand{\safeimage}[2]{%
    \IfFileExists{#1}{\includegraphics[width=#2\textwidth]{#1}}{%
    \fbox{\parbox[c][0.18\textheight][c]{#2\textwidth}{\centering Missing generated figure\\\codepath{#1}}}}}

\title{{\Large EP-392: UG Project-2}\\[1.2em]
Downscaling Approach To Produce High Spatial Soil Moisture Products Using Satellite Data}
\author{Mayur Kunnitigi Matada - 23174017\\[0.5em]
Supervisor: Prof. Rajendra Prasad\\Department Of Physics}
\date{October 2026\\[2.0em]\includegraphics[width=0.55\textwidth]{figures/IIT-BHU_Logo.pdf}}

\begin{document}
\maketitle
\thispagestyle{empty}
\clearpage
\tableofcontents
\clearpage

\chapter*{From Regional Foundations to Global Scale}
\addcontentsline{toc}{chapter}{From Regional Foundations to Global Scale}

Stage One established the regional foundation for this project through a focused
study of soil-moisture downscaling over Spain. The work used REMEDHUS station
observations as the reference target and combined them with satellite,
meteorological, soil, and terrain predictors. It first developed a quality
controlled data pipeline and a residual XGBoost baseline, and then used the
same carefully checked inputs to build a multimodal spatial-temporal deep
model. Evaluation was performed under an unseen-station protocol so that the
reported improvement measured transfer to stations not used for training,
rather than simply memorisation of randomly held-out rows.

The main outcome of Stage One was therefore more than a single accuracy value. It
established the data-cleaning rules, daily target construction, patch-based
feature representation, residual anchoring strategy, and station-disjoint
evaluation approach that support the wider project. It also showed why
neighbourhood context, temporal information, and explicit handling of the
coarse soil-moisture anchor are useful when a point observation is compared
with area-integrated satellite and meteorological measurements.

Stage One identified globalization as the next major step. A regional result over
Spain cannot by itself demonstrate that the method transfers across climate
zones, soil types, vegetation regimes, station networks, and satellite
coverage conditions. Stage Two addresses that future-work objective by extending
the system to a global, multi-provider dataset and by redesigning the
acquisition, storage, alignment, caching, training, and audit machinery for
millions of station-date examples. The global system retains the scientific
lessons of Stage One while making the computational assumptions explicit enough
to support a much larger and more diverse experiment.

\chapter*{Abstract}
\addcontentsline{toc}{chapter}{Abstract}

This report documents the construction of a global, multi-modal soil-moisture
learning system from raw observation archives through remote data acquisition,
storage engineering, model training, calibration, and scientific error analysis.
The project is substantially larger than a conventional machine-learning
experiment: its reported implementation footprint is 143,850 lines of code,
including archive parsers, scientific data contracts, cloud-raster clients,
distributed-work recovery, array conversion, validation, model training, and
diagnostic analysis.

The target is near-surface volumetric soil moisture in \si{m^3/m^3}. The
reference observations come from the International Soil Moisture Network (ISMN)
and are quality filtered, depth classified, and aggregated to UTC calendar-day
targets. Predictors combine Copernicus ERA5-Land meteorological forcing, NASA
SMAP microwave context, MODIS and Sentinel-1 products hosted by Microsoft
Planetary Computer, and static soil and terrain layers extracted through Google
Earth Engine. The provenance distinction is essential: Google Earth Engine
supplies static layers only; Planetary Computer remains the authoritative host
for the dynamic raster products; and a paid Azure Poland virtual machine was
used as the high-throughput processing client after network optimisation.

The final v19.5 model uses modality-specific residual spatial encoders, global
mean/max/attention pooling, a temporal convolutional network with irregular
history masking and last/history gating, static and location branches, a learned
five-way modality gate, nonlinear fusion, and an identity-initialised quadratic
calibration layer. The location-history configuration reaches calibrated test
RMSE $0.06594$ and $R^2=0.65341$ over 905,248 rows; the environment-only
configuration reaches RMSE $0.08626$ and $R^2=0.47046$ over 644,716 rows.

The report emphasizes technical feats and their scientific boundaries. The
system reduces a naive estimate of roughly 700~GB of data and 60 days of
top-tier GPU computation through scene clustering, COG overview reads, resumable
shards, chunk-level caching, memory-mapped arrays, train-only FP16
normalisation, bounded prefetch, gradient accumulation, and staged
hyperparameter search. These are not merely performance tricks: they determine
whether the global experiment can be repeated, audited, and interpreted.

\chapter*{Executive Summary}
\addcontentsline{toc}{chapter}{Executive Summary}

\section*{Project Overview}
The project transformed heterogeneous, irregular, multi-provider global
observations into a verified and resumable scientific-computing pipeline that
can train and evaluate a multi-modal neural model at a scale that would be
impractical under naive data and compute assumptions.

\section*{Quantified scale}
\begin{itemize}
    \item 143,850 reported lines of implementation.
    \item 4,340 ISMN station files across 61 networks.
    \item 144,466,562 observations audited.
    \item 103,504,765 good-quality observations retained in the audit.
    \item 5,045,849 daily Parquet rows and 3,514,643 unique station-date samples.
    \item 204,033 unique Sentinel-1 scenes indexed.
    \item 528,280 scene-cluster associations recorded.
    \item 21,655 resumable NPZ scene shards produced.
    \item 905,248 held-out location-history predictions evaluated.
\end{itemize}

\section*{What makes the work technically difficult}
The data sources do not share a time grid, spatial grid, resolution, missingness
pattern, quality convention, storage format, or network access mechanism. A
station observation is a point-like measurement at a particular depth; a
satellite product is an area-integrated measurement whose acquisition date may
be offset from the target; reanalysis is a gridded forcing field; static layers
describe slow site properties; and a microwave product such as SMAP has a large
footprint. The model is therefore only one stage of the scientific problem.

\feat{The project solved data identity before model identity. A feature array
with the right shape can still be scientifically wrong if its row order has
drifted. Stable sample identifiers, producer row indexes, sidecar tables,
completion flags, shape checks, sampled reads, and checksums were used to make
association auditable.}

\feat{The project solved remote request geometry before simply adding more
workers. STAC searches and COG reads were clustered by shared scene and time
window, overview reads were benchmarked against the 1~km target, and finished
scene units were persisted so failures did not force a global restart.}

\feat{The project solved model throughput before treating GPU size as the only
limitation. Normalised FP16 memory maps, vectorised gathers, bounded queues,
pinned memory, asynchronous copies, compilation fallback, microbatches,
gradient accumulation, and EMA reduced the wall-clock cost and the failure
rate of training.}

\section*{Summarized Results}
\begin{center}
\begin{tabular}{lrrrrr}
\toprule
Configuration & Test rows & RMSE & MAE & $R^2$ & Robust $R^2$\\
\midrule
Environment-only & 644,716 & 0.08626 & 0.06794 & 0.47046 & 0.51276\\
Location-history & 905,248 & \textbf{0.06594} & \textbf{0.04854} &
\textbf{0.65341} & \textbf{0.74078}\\
\bottomrule
\end{tabular}
\end{center}

\feat{The location-history result is a calibrated result. Its raw test RMSE
is approximately 0.11190 before the learned output calibration. The report
therefore distinguishes raw, calibrated, robust, and anchor comparisons rather
than presenting one score as if it represented every layer of the system.}

\chapter{Project Motivation, Scope, and Scientific Question}

\section{Soil moisture as a physical state variable}
Soil moisture is the water held in soil pore space and is commonly represented
as volumetric water content. In the surface layer it responds to precipitation,
infiltration, evaporation, transpiration, runoff, soil texture, terrain,
vegetation, and atmospheric demand. It is not a passive scalar that changes
uniformly across a landscape. Two neighbouring locations can respond
differently to the same rainfall because their storage and drainage properties
differ.

The physical importance is broad. Surface moisture influences land-atmosphere
energy exchange, plant stress, drought evolution, runoff generation, and
irrigation demand. The project is therefore motivated by a measurement
problem: in-situ stations provide high-value local observations but are sparse,
whereas remotely sensed products provide spatial coverage but are indirect,
coarse, intermittent, or affected by clouds, vegetation, geometry, and
footprint mismatch.

\section{Why downscaling and fusion are needed}
A coarse microwave retrieval can provide a useful large-scale state estimate
while averaging over fields, roads, terrain, and vegetation communities. A
finer-scale product must use auxiliary information to infer local variation.
This is not the same as creating new measurements. It is a conditional
reconstruction constrained by available observations and learned relationships.

The global system combines:
\begin{enumerate}
    \item An in-situ target with explicit quality and depth policy;
    \item Optical vegetation and thermal information from MODIS;
    \item Radar backscatter structure from Sentinel-1;
    \item Meteorological forcing and memory from ERA5-Land;
    \item A coarse microwave anchor from SMAP;
    \item Static texture and terrain from GEE exports;
    \item Location and station history in one controlled model configuration.
\end{enumerate}

\section{The central research question}
The operational research question is:
\begin{quote}
Can a verified multi-modal system estimate daily near-surface soil moisture at
1~km station-centred support from heterogeneous global observations while
retaining enough provenance, computational efficiency, and diagnostic
transparency for scientific interpretation?
\end{quote}

This question contains three subquestions:
\begin{enumerate}
    \item Can the data sources be harmonised without silently changing the
    physical target?
    \item Can the remote-sensing workload be executed at global scale without
    repeated or unrecoverable work?
    \item Does a model with spatial, temporal, static, location, and anchor
    pathways outperform simpler environmental-only alternatives under
    station/network-aware evaluation?
\end{enumerate}

\section{Scope and non-claims}
The model estimates the observed target distribution; it does not solve a
complete soil-water balance, perform causal inference, or guarantee physical
conservation. The report does not claim that a high $R^2$ proves universal
geographic transfer. It also does not treat oracle trimming as a deployment
filter or confuse the Azure processing client with the Planetary Computer
archive.

\chapter{Architecture of the Entire Scientific-Computing System}

\section{System-level view}
\begin{figure}[H]
\centering
\includesvg[inkscapelatex=false,width=0.65\textwidth]{Builtsvg/latex/1}
\caption{The global system is a chain of scientific contracts and computing systems, not only a neural-network block. Each stage reduces ambiguity before the next stage increases computational complexity.}
\label{fig:system}
\end{figure}

\begin{landscape}
\begin{figure}[H]
\centering
\includesvg[inkscapelatex=false,width=\linewidth,height=0.78\textheight,keepaspectratio]{Builtsvg/latex/2}
\caption{Full expanded production DAG. The diagram deliberately exposes the intermediate operations that are normally hidden inside broad pipeline boxes.}
\label{fig:full-dag}
\end{figure}
\end{landscape}

\section{Scale of the project}
The 143,850-line implementation is distributed across executable notebooks,
Python modules, shell orchestration, validation programs, model exports, and
analysis outputs. The number is useful as a description of scope rather than
as a quality score: the repository contains separate systems for scientific
contracts, observation parsing, gridded environmental retrieval, STAC/COG
access, distributed recovery, storage conversion, cache construction, model
training, and evaluation.

The scale is visible in the work products. The project audited 144 million
observations, reduced them to more than 3.5 million unique station-date
identities, indexed 204,033 Sentinel-1 scenes, associated those scenes with
528,280 cluster records, and produced 21,655 resumable shards. The final
evaluation then operated over hundreds of thousands of held-out predictions.
Each count corresponds to a different systems problem: streaming and audit,
identity preservation, remote request amortisation, restartability, and
full-population analysis.

\chapter{Global Data Contract and Ground Truth}

\section{Why a contract is necessary}
When data arrive from dozens of networks and products, a model can train
successfully on an invalid mixture. A contract makes the physical meaning
explicit before storage and model code. The current global contract records
UTC dates, a nominal 1~km grid, a $13\times13$ footprint, target units,
accepted quality, depth classes, and channel counts.

\section{Depth and quality}
The target policy distinguishes near-surface 0--5~cm measurements from the
explicit 5--10~cm fallback and excludes measurements outside 0--10~cm. Only
quality flag \texttt{G} contributes to daily labels. This matters because a
numerically finite value from a deeper sensor is not interchangeable with a
surface value.

\section{Archive audit and streaming conversion}
The archive contains 4,340 station files, 4,002 compatible-depth files, and
61 networks. The converter supports directories and ZIP archives without
extracting the complete archive. One source file is opened, its header is
parsed, valid observations are accumulated by UTC day, and rows are written
to Zstandard-compressed Parquet.

\begin{lstlisting}[style=pycode, caption={Conceptual streaming reduction used by the global converter}]
for source_file in iter_source_files(source):
    header = parse_header(source_file)
    if not compatible_depth(header):
        continue
    by_day = defaultdict(list)
    for fields in source_file.observations:
        if fields.quality != "G":
            continue
        day, value = parse_finite_measurement(fields)
        by_day[day].append(value)
    rows = daily_means_with_provenance(header, by_day)
    parquet_writer.write_table(rows)
\end{lstlisting}

\feat{The achievement is not just compression. Streaming changes the maximum
working set from “the entire archive plus dataframe overhead” to “one source
file plus one daily accumulator.” The resulting Parquet table remains
columnar, compressed, auditable, and joinable.}

\section{Ground-truth statistics}
\begin{itemize}
    \item 144,466,562 observations audited.
    \item 103,504,765 good-quality rows.
    \item 40,961,797 non-good rows.
    \item Zero malformed rows in the reported global profile.
    \item 5,045,849 daily Parquet rows.
    \item 2,210 stations and 61 networks in the daily profile.
    \item 3,514,643 unique station-date samples.
\end{itemize}

\feat{The headline benchmark combines primary and fallback depth classes.
Depth-stratified reporting remains important because the hydrological response
time and storage physics are not identical between 0--5 and 5--10~cm.}

\chapter{Data Sources, Provenance, and Environmental Predictors}

The predictors are introduced in the order in which they contribute to the
physical interpretation of the target. Meteorological forcing describes the
conditions that drive change, microwave observations provide a direct but
coarse moisture-related anchor, optical and radar products add local spatial
context, and static layers describe persistent differences between locations.
Keeping these roles distinct helps explain both the value of fusion and the
limits of treating every input as interchangeable.

\section{Data sources and providers}
The global dataset combines observations and environmental products produced
or maintained by different scientific organisations. Keeping the provider
explicit is important because the organisation that hosts a file is not
always the organisation that created the measurement or product.

\begin{itemize}
    \item \textbf{International Soil Moisture Network (ISMN):} in-situ soil-moisture observations and quality flags. The ISMN is maintained as an international network, with measurements contributed by individual monitoring networks and station operators.
    \item \textbf{ERA5-Land:} meteorological reanalysis, including precipitation and temperature fields. It is produced by the Copernicus Climate Change Service and the European Centre for Medium-Range Weather Forecasts (ECMWF).
    \item \textbf{Soil Moisture Active Passive (SMAP):} coarse-resolution microwave soil-moisture observations and the Level-4 soil-moisture product. SMAP is a NASA mission and its products are produced by NASA.
    \item \textbf{MODIS:} vegetation and land-surface-temperature products. MODIS is an instrument operated by NASA aboard the Terra and Aqua satellites; the products are produced by NASA.
    \item \textbf{Sentinel-1 RTC:} radar backscatter products used for VV and VH features. Sentinel-1 is an Earth-observation mission of the European Space Agency (ESA). The cloud-hosted assets used here are accessed through Microsoft Planetary Computer.
    \item \textbf{Microsoft Planetary Computer:} the Microsoft-hosted STAC catalogue and cloud storage interface used to discover and retrieve MODIS and Sentinel-1 assets. It is an access and hosting platform, not the scientific producer of those satellite missions.
    \item \textbf{Google Earth Engine:} the Google geospatial-computing platform used to extract the static soil and terrain layers. The platform provides the processing environment; the underlying products retain their original producers.
    \item \textbf{SoilGrids:} global soil-property maps used for clay, sand, and silt information. SoilGrids is produced by ISRIC---World Soil Information.
    \item \textbf{NASADEM:} elevation data used to derive terrain variables such as slope and aspect. NASADEM is produced from NASA and partner satellite radar data and distributed through NASA's Earth-observation data systems.
\end{itemize}

\begin{figure}[H]
\centering
\includegraphics[width=\textwidth]{figures/data_sources/ISMN_global_station_map.png}
\caption{Global distribution of stations in the International Soil Moisture Network (ISMN), showing the worldwide observation network from which the global in-situ data are drawn. The map is reproduced as a reference to the ISMN network and the Austrian GCOS inventory report \citep{GCOSAustria}.}
\label{fig:ismn-global-stations}
\end{figure}

\section{Provenance map}
\begin{figure}[H]
\centering
\includesvg[inkscapelatex=false,width=\textwidth,height=0.72\textheight,keepaspectratio]{Builtsvg/latex/3}
\caption{Data provenance and infrastructure roles. Azure accelerates processing; it does not replace Planetary Computer as the data host.}
\label{fig:provenance}
\end{figure}

\section{Copernicus ERA5-Land}
ERA5-Land supplies daily meteorological forcing. The feature layer includes
precipitation, 2-m temperature, dewpoint, vapour-pressure deficit, and an
antecedent precipitation index. The access strategy reuses unique reanalysis
grid points and processes point batches rather than repeatedly requesting the
same grid location for nearby stations.

\section{NASA SMAP}
SMAP provides a broad-footprint microwave soil-moisture context. The extraction
aggregates valid observations to daily station or cell values and retains
observation counts, conflicts, and availability. Missing SMAP is represented
explicitly rather than filled with a physical-looking zero.

\section{MODIS from Planetary Computer}
MODIS supplies EVI and LST through STAC discovery and COG or tile access.
Year-partitioned Zarr arrays preserve the multidimensional product while keeping
the scalar sample index separate. The EVI and LST arrays retain their distinct
footprints rather than being forced into a false common native resolution.

\section{Sentinel-1 RTC from Planetary Computer}
Sentinel-1 provides radar-sensitive VV and VH context and their ratio. The
nearest valid scene is selected within an 18-day tolerance, with deterministic
earlier-date tie breaking. The production path uses a benchmarked 512 by 512
overview to approximate the 1~km target while retaining a separate exact
native-resolution route.

\feat{Remote-sensing extraction is request-level scientific engineering. STAC
searches, SAS signing, COG windows, overview semantics, acquisition age,
cluster mapping, shard completion, and checksums all contribute to the numerical
meaning of the final feature.}

\section{Static layers from Google Earth Engine}
GEE contributes static soil and terrain variables: clay, sand, silt, elevation,
slope, and aspect. Static extraction is separated from dynamic acquisition and
reused across years. Older GEE notebooks are treated as experiments or
prototypes, not as evidence that dynamic products came from GEE.

\chapter{Remote Acquisition and Distributed Data Engineering}

Once the labels and predictor roles were defined, the next challenge was to
obtain the data without losing either provenance or operational control. The
remote products are large, intermittently accessed, and hosted through
different interfaces. The following sections therefore describe not only where
the data came from, but also how the extraction was organised so that a
temporary network or machine failure did not become a scientific failure.

\section{Why remote data extraction became the dominant systems problem}
Remote raster extraction has fixed costs beyond pixel count: catalogue search,
authentication and signing, TLS and HTTP setup, COG tile lookup, decompression,
window conversion, reprojection, and local serialization. A design that performs
one complete sequence for every station-date sample can be dominated by fixed
overhead.

\section{Cell reduction and collision maps}
Stations are mapped to deterministic approximate 1~km cells. Collocated requests
are collapsed for the expensive raster operation while a collision map retains
the relationship between original station IDs and the representative cell.
This is a major reduction in redundant work, but the collision map prevents the
reduction from becoming an undocumented loss of identity.

\section{Scene clustering}
\begin{figure}[H]
\centering
\includesvg[inkscapelatex=false,width=\textwidth,height=0.72\textheight,keepaspectratio]{Builtsvg/latex/4}
\caption{The clustering transformation changes the unit of work from a sample to a reusable spatial/temporal cluster.}
\end{figure}

The audited Sentinel state contained 819 cluster searches, 204,033 unique
scenes, 528,280 scene-cluster associations, 818 clusters with associations,
21,655 NPZ shards, and no unlinked scenes. These are not merely bookkeeping
numbers. They demonstrate that the remote archive was represented as a
recoverable graph of scenes, clusters, and samples.

\section{Azure Poland processing client}
The Planetary Computer archive remained authoritative. After network-read
optimisation, a paid Azure VM located in Poland served as the high-throughput
client. The VM supplied a better-connected execution environment and adequate
disk and memory for concurrent COG work. It was not a replacement archive.

\section{Resume engineering}
The workflow uses SQLite backup operations, cluster search state, scene shards,
completion markers, manifests, and SHA-256 checksums. A failure therefore
invalidates incomplete units rather than erasing all completed work.

\feat{This is analogous to checkpointing a long numerical simulation. A global
remote extraction is decomposed into verified subproblems, and the computation
can resume from the last trusted state. The same principle is later applied to
local caches and model checkpoints.}

\section{Overview versus exact reads}
The 1~km target does not necessarily require repeatedly downloading every
10-metre Sentinel-1 pixel. Overview reads reduce the remote payload and
decompression cost. Exact reads remain available for validation and sensitivity
analysis. Separate namespaces prevent approximate and exact arrays from sharing
a filename and silently changing scientific meaning.

\chapter{Storage Formats, Merge, and Alignment}

The extraction stage produces several kinds of information: scalar tables,
mutable job state, multidimensional rasters, and completed evidence products.
These objects have different access patterns, so the storage design is part of
the scientific method rather than an implementation detail. The sections below
explain how each representation is used and how the joins between them are
checked.

\section{Why no single file format is sufficient}
The project uses each representation for a distinct access pattern:
\begin{itemize}
    \item ZIP streaming for the compressed observation archive;
    \item Parquet for compressed, columnar scalar tables;
    \item SQLite for indexed mutable scene state;
    \item NPZ for completed transferable scene shards;
    \item Zarr for chunked multidimensional arrays;
    \item NumPy memory maps for local random access during training;
    \item JSON and CSV for human-readable manifests and provenance.
\end{itemize}

Choosing a universal format would force a compromise between streaming text,
indexed state, cloud rasters, and GPU batch reads. The storage architecture is
therefore an important part of the scientific design.

\section{Low-memory merge}
The merge path retains raster arrays in yearly modality-specific Zarr stores
and joins scalar tables through stable sample identities and explicit row-index
columns. DuckDB and Parquet scans avoid loading the complete global table into
one pandas object.

\section{The row-order hazard}
A row-order error is more dangerous than a visible exception. If a Sentinel
array has the correct shape but its sidecar IDs are in a different producer
order, the model can train on plausible but incorrect pairs. The alignment
checks therefore inspect row counts, sidecar IDs, producer sort contracts,
expected patch shapes, metadata, completion flags, and sampled values.

\begin{lstlisting}[style=pycode, caption={Conceptual row-association validation}]
assert len(sidecar) == zarr_array.shape[0]
assert sidecar["sample_id"].is_unique
assert np.array_equal(sidecar["row_index"], np.arange(len(sidecar)))
assert finite(sampled_rows)
assert expected_patch_shape(zarr_array)
\end{lstlisting}

\feat{Index alignment establishes association, not geophysical equivalence.
CRS, affine transform, pixel-centre offsets, overview resampling, and physical
support area still require modality-specific geolocation checks.}

\chapter{The Optimization Programme}

The optimisation programme should be read as a sequence of responses to
specific failure modes. The aim was not to make every individual operation
faster in isolation. Instead, the design reduced unnecessary work, moved
repeated work to an earlier stage, kept the CPU and GPU supplied at a
controlled rate, and made failure recoverable. This distinction matters
because a global experiment can be slow even when its arithmetic is efficient:
the dominant cost may be waiting for data, moving data, or repeating work that
has already been completed.

\section{The counterfactual}
The counterfactual makes the need for these choices concrete. Without
optimisation, repeated remote reads, expanded patch tables, incompatible
intermediates, per-batch raster decoding, and GPU starvation were estimated to
produce approximately 700~GB of data and 60 days of top-tier GPU compute. The
network transfer requirement would also be impractical. More importantly, this
would have been a fragile calculation: a failure late in the run could have
forced large portions of the earlier work to be repeated.

\section{Optimization principle one: do not create data that will be discarded}
The first decision was to prevent unnecessary data from entering the expensive
part of the pipeline. Quality and depth filtering occur before expensive joins,
so rejected observations do not trigger remote extraction. Daily labels replace
redundant intraday rows for the model target while preserving their provenance.
ERA5 requests are reused at unique grid points, static layers are exported
once, and label-date manifests restrict dynamic extraction to dates for which a
target exists. Finally, only the active Sentinel-1 bands enter the production
model cache. Each of these choices is small by itself, but together they reduce
the number of rows, dates, bytes, and conversions that every later stage must
handle.

\section{Optimization principle two: amortize remote fixed costs}
Remote operations have a fixed cost before any useful pixel is returned:
connection setup, authentication, catalogue search, asset opening, and block
decompression. Scene clustering, cached STAC searches, COG-aligned windows,
overview reads, and shared scene opens amortise those costs over many
station-date rows. The Azure location then makes the resulting request pattern
practical by reducing the network distance and improving sustained transfer.

\section{Optimization principle three: cache at the storage grain}
The most useful cache unit is the unit the storage system can retrieve
efficiently. Individual patch caches had little reuse because the
unique/requested index ratio was close to one: many requested patches touched
different source locations. The system instead caches native source chunks,
decodes each touched chunk once, normalises it once, and gathers requested rows
in batch order. This avoids paying the decompression and conversion cost again
for every training example.

\section{Optimization principle four: move deterministic work off the GPU}
The GPU should spend its limited time on operations that benefit from massive
parallel arithmetic. The local normalised cache therefore performs fill
handling, physical scale conversion, validity detection, training-only
standardisation, clipping, and FP16 conversion once. Training can then read
value and mask channels from memory-mapped arrays rather than repeatedly asking
the GPU to perform the same preparation for every epoch.

\section{Optimization principle five: bound parallelism}
Parallelism is useful only while the shared resources can support it. Unbounded
workers and prefetch queues can exhaust RAM, thrash the disk cache, or trigger
remote throttling. The training system therefore caps workers, prefetch depth,
and future-batch queues. Pinned memory and asynchronous CUDA copies are used
within that bounded system so that short delays can be hidden without turning
the input pipeline into an uncontrolled memory consumer.

\section{Optimization principle six: preserve the effective batch under finite VRAM}
The desired effective batch was 2,048, but the available VRAM could not hold all
of those examples and their intermediate activations in one allocation.
Microbatches and gradient accumulation preserve the effective batch without
requiring that single allocation. Automatic mixed precision reduces the memory
occupied by suitable intermediate values, while gradient clipping, EMA, and
best-checkpoint restoration improve stability when a long run encounters
unusually difficult batches or noisy epochs.

\section{Optimization principle seven: make caches self-describing}
An optimisation is dangerous if it makes stale data look current. Scaler hashes
and build signatures prevent arrays generated under different decode rules,
years, or normalisation settings from being mixed. Temporary
\texttt{.building} files are promoted only after validation, so an interrupted
construction cannot silently replace the last complete cache.

\section{Optimization register in prose}
The engineering decisions are interdependent. Streaming ingestion reduces host
memory; Parquet reduces persistent scalar volume; clustered extraction reduces
remote fixed costs; Zarr preserves spatial arrays; chunk caching reduces
decompression; normalised memory maps reduce CPU work; bounded prefetch reduces
GPU idle time; and gradient accumulation preserves optimisation behaviour under
VRAM limits. Treating these as isolated tricks would miss their multiplicative
effect.

\section{What the optimisation means in computer terms}
For a physics reader, the central distinction is between computation and
movement. A processor may be able to calculate a mean over a raster window very
quickly, but it cannot calculate until the bytes have arrived from the remote
server, been authenticated, decompressed, copied into memory, and arranged in
the expected array. In the original station-date formulation, those fixed
steps were repeated for many rows that shared the same scene or grid cell.
Clustering changes the accounting: one expensive remote read produces a
reusable object, and many scientific rows point to that object through an
index.

The network optimisation therefore works at three levels. First, the catalogue
search is cached so the same question about available scenes is not sent again.
Second, nearby windows are grouped so one connection and one scene open can
serve several requests. Third, COG overviews select a lower-resolution internal
tile when the target is approximately one kilometre, avoiding transfer of
unneeded ten-metre pixels. The third step is not free: it changes the support of
the measurement. That is why the exact path is retained, the strategy is named
in the manifest, and exact-versus-overview comparisons remain a required
validation.

The storage optimisation addresses a different resource: random access. A
compressed archive is efficient for distribution but poor for repeatedly
extracting arbitrary rows. A full dense array is easy to index but can exceed
RAM and takes too long to rebuild after a failure. Chunked Zarr arrays reduce
the amount read from disk; memory-mapped arrays let the operating system bring
only requested pages into memory; and precomputed normalisation means the
training loop does not repeatedly perform the same conversion. The validity
mask is stored beside the value because replacing an absent value with a
numeric fill would make missingness indistinguishable from a physical state.

The CPU and GPU optimisations address a pipeline-balance problem. If the CPU
prepares batches more slowly than the GPU consumes them, the GPU waits. If
workers prepare too far ahead, RAM fills with batches that may never be used.
Bounded worker counts and prefetch queues keep the queue between those two
failure modes. Pinned memory and asynchronous copies overlap host-to-device
transfer with arithmetic where possible. Mixed precision reduces the size of
activations, while accumulation allows many small device batches to produce the
gradient statistics of a larger effective batch. These changes do not mean the
model performs fewer conceptual updates; they change when bytes move and when
the optimiser commits an update.

Finally, restartability is a time optimisation as much as a reliability
feature. A 700~GB-style run is not one indivisible calculation; it is a
collection of searches, scenes, chunks, shards, and manifests. A completion
marker is a promise that a particular unit passed its checks. After an
interruption, the scheduler can spend time only on units without a valid
promise. Checksums prevent a truncated file from being mistaken for a completed
unit, and cache signatures prevent an old unit produced under different
settings from being reused. The result is fewer wasted hours, fewer repeated
network transfers, and a more interpretable scientific run.

\feat{The value of the optimisation programme is not a single speedup. It is
that an experiment which could fail for storage, network, memory, or restart
reasons became an experiment whose assumptions are explicit and whose
intermediate states can be verified.}

\section{Why the unoptimised experiment was not practical}
It is useful to describe the original difficulty as a queue of physical
operations rather than as an abstract ``large dataset''. A request began at a
station and date, contacted a catalogue service, obtained an authorised asset
URL, opened a remote raster, downloaded compressed blocks, decompressed them,
converted their numeric encoding to physical units, handled missing pixels,
assembled a patch, and finally wrote the result. Only after all of those steps
could one small training row be produced.

The same scene, grid cell, or source chunk was often needed by many rows. If
each row repeated the complete sequence, the experiment paid the fixed costs
again and again. The expensive part was therefore not only the arithmetic
needed for the final mean or neural-network prediction. It was the movement of
data: network latency, authentication, transfer, decompression, disk reads,
memory copies, and conversion between representations. A faster processor
cannot remove time spent waiting for a remote server or a disk.

\subsection{A simple analogy for the processing pipeline}
The computer can be viewed as a factory with four work areas. The remote
archive is the supplier, the CPU is the preparation area, the GPU is the
specialised high-speed machine, and the disk is the warehouse. The GPU is very
fast at tensor arithmetic but has limited local memory. If the supplier sends
parts one at a time, the GPU stands idle. If preparation sends parts faster
than the GPU can consume them, the warehouse fills and the machine may run out
of memory. A successful design must keep all four areas busy without allowing
any one queue to become unbounded.

This analogy also explains why several modest optimisations were needed
together. Clustering improved the supplier connection, chunk caching reduced
warehouse retrieval, precomputed normalisation reduced preparation, bounded
prefetch controlled the queue, and microbatching adapted the work to the
GPU's limited local memory. None of these changes alone solves the complete
problem.

\subsection{Network traffic: reduce requests before reducing workers}
Adding more workers to a poorly organised remote workload would have made the
problem worse. Every worker would have repeated catalogue searches, opened the
same scenes, and competed for the same network and remote service limits. The
pipeline first changed the unit of work. Station dates were grouped into
canonical cells and spatial/temporal scene clusters. One search result and one
opened source scene could then serve many station-date rows.

The STAC search cache records the answer to a catalogue question. The scene
inventory records which source scenes can serve which clusters. Signed asset
URLs and request metadata are retained so a retry does not have to rediscover
the entire remote state. This changes the workload from ``ask the archive for
every row'' to ``discover the reusable source once, then reference it by
index''. The scientific identity remains the station-date row; only the
computer's route to the data changes.

\subsection{Why overview reads helped, and what they changed}
Sentinel-1 native pixels are much smaller than the approximately one-kilometre
model target. Reading every native pixel for every sample would transfer and
decompress information that would later be averaged away. Cloud-Optimised
GeoTIFF overviews provide precomputed lower-resolution representations, so the
remote read can request a smaller block when that support is scientifically
appropriate.

This is a performance approximation, not a free replacement for the source.
An overview represents a different spatial support and may use a different
resampling path. The extraction therefore records whether an overview or exact
native read was used, keeps exact reads available for validation, and prevents
the two strategies from sharing an ambiguous filename. The optimisation saves
network bytes and decompression work while preserving the information needed to
quantify the approximation.

\subsection{Memory: why a correct calculation can still crash}
Memory is not one single pool. The GPU has fast but limited VRAM; the CPU has
larger system RAM; and the operating system uses disk-backed pages when RAM is
under pressure. A batch can fail even when the final model is small because
the batch also requires intermediate activations, gradients, temporary arrays,
worker queues, and copies between host and device.

The production cache avoids constructing the entire global tensor in RAM.
Zarr stores multidimensional arrays in chunks, so a request loads the touched
chunks rather than the whole year. Memory maps expose large prepared arrays as
file-backed regions; the operating system brings in the pages actually used by
the next batch. A value channel is stored beside a validity channel, so a
missing pixel is not confused with a genuine physical zero or fill value.

Normalisation, clipping, fill handling, and conversion to the training
precision are performed once while the cache is built. Without this step, the
same deterministic work would be repeated for every epoch. The training loop
then performs the work that must change sampling, forward prediction,
backpropagation, and optimisation instead of repeatedly decoding the source
archive.

\subsection{GPU memory: microbatches preserve the learning problem}
The desired optimisation batch was 2,048 examples, but placing 2,048 complete
examples and all their intermediate tensors on the GPU at once would exceed
VRAM. The solution was not to change the intended batch arbitrarily. A
microbatch that fits in VRAM is processed, its gradients are accumulated, and
several microbatches are combined before one optimiser update. Thus the
physical batch controls memory while the effective batch preserves the
intended update scale.

Automatic mixed precision stores many intermediate values in a smaller
floating-point representation when safe, reducing memory traffic and often
accelerating tensor operations. Gradient clipping limits the effect of an
exceptionally large update. Exponential moving averages provide a smoother
parameter state for checkpoint selection. These mechanisms address different
failure modes: VRAM pressure, numerical instability, and sensitivity to the
last few noisy updates.

\subsection{CPU GPU balance: keeping the expensive machine occupied}
The GPU cannot start a batch until the CPU has gathered the requested rows,
assembled temporal history, applied masks, and copied the result to device
memory. If this preparation takes longer than the GPU computation, utilisation
falls even though the GPU itself is capable of much more work. Conversely,
unlimited workers can prepare too many batches, consume RAM, and cause
contention that makes every worker slower.

The data loader therefore uses a bounded number of workers and a bounded
prefetch queue. Pinned host memory makes transfers more predictable, and
asynchronous copies can overlap transfer with computation where the hardware
allows it. The queue is deliberately bounded: it is large enough to hide short
storage delays but small enough that an unused future batch cannot exhaust
memory. This is a control problem, not simply a request to maximise the worker
count.

\subsection{Storage formats: each format removes a different bottleneck}
The project uses several formats because the operations are different. ZIP is
useful for distributing the original archive but poor for random row access.
Parquet is efficient for columnar tables and low-memory joins. SQLite provides
indexed mutable state for scene searches and completion records. Zarr supports
chunked raster arrays. NPZ shards make completed remote units portable and
independently checkable. Memory-mapped NumPy arrays provide fixed-offset
training access.

Converting everything into one format would move the bottleneck rather than
remove it. The important engineering rule is that every conversion preserves a
mapping back to the canonical sample identity. Shapes alone are insufficient:
an array with the right dimensions can still contain the wrong rows. Sidecar
row indexes, manifests, finite-value checks, completion markers, and checksums
make the cheaper representation auditable.

\subsection{Restartability: saving time when failure is expected}
At global scale, interruption is a normal operating condition. A network
timeout, expired asset, worker failure, full disk, or machine restart should
not force the project to repeat all previous work. The extraction is therefore
divided into independently committed units such as catalogue searches, scene
clusters, source chunks, and shards.

An output is not considered complete merely because a filename exists. The
workflow checks dimensions, source identifiers, valid-value counts, metadata,
and checksums before recording completion. Writes use temporary paths and are
promoted atomically only after validation. On restart, the scheduler trusts
only units with a valid completion record and retries the rest. This converts a
potentially all-or-nothing computation into a sequence of smaller verified
computations.

\subsection{The resulting saving in practical terms}
The optimisations saved more than wall-clock time. They reduced repeated
network transfer, reduced remote catalogue pressure, reduced peak RAM and VRAM,
reduced disk churn, reduced GPU idle time, and reduced the amount of work lost
after an interruption. They also made the result explainable: a scientist can
identify which source asset, read strategy, cache signature, row mapping, and
checkpoint produced a result.

\begin{center}
\begin{tabular}{p{0.25\textwidth}p{0.32\textwidth}p{0.31\textwidth}}
\toprule
\textbf{Bottleneck} & \textbf{Naive behaviour} & \textbf{Production response}\\
\midrule
Remote archive & Repeat searches, authentication, scene opens, and downloads
for individual rows & Cell reduction, scene clustering, cached searches,
shared scene reads, and explicit overview strategy\\
CPU and disk & Repeated decompression, conversion, and random reads during
every epoch & Chunked Zarr, source-chunk cache, normalised mmap arrays, and
vectorised gathers\\
GPU memory & Attempt to place the full optimisation batch on limited VRAM &
Microbatches, gradient accumulation, AMP, and batch-size probing\\
Host memory & Unlimited workers and prefetch create an unbounded queue &
Bounded workers, bounded prefetch, pinned memory, and finite caches\\
Failures & One interruption invalidates a long monolithic run & Atomic shards,
completion markers, manifests, checksums, and resume logic\\
\bottomrule
\end{tabular}
\end{center}

\chapter{Engineering Internals of the Global Run}

The preceding chapters described the pipeline from the outside. This chapter
now follows the data through the computer itself: which unit is processed at
each stage, how intermediate representations are stored, how a failed unit is
distinguished from a completed one, and how the training loop moves data
between storage, CPU memory, and GPU memory. This view connects the engineering
decisions to the scientific guarantees reported later.

\section{The unit of computation was redesigned several times}
The first conceptual unit of the experiment was a station-date sample. That
unit is scientifically natural because the target is defined at a station and
date, but it is computationally inefficient. Many station dates fall in the
same raster cell, request the same reanalysis grid point, or refer to the same
remote scene. Treating every row as an independent download would repeat
network requests, authentication, decompression, and array conversion.

The production system therefore used different units at different stages. A
station-date remained the unit of scientific identity. A canonical cell became
the unit of spatial request sharing. A scene-cluster became the unit of remote
Sentinel work. A yearly Zarr chunk became the unit of persistent raster storage.
A packed tensor row became the unit of training. This separation is important:
it reduces repeated work without changing the definition of the target row.
Each conversion writes an index that maps the cheaper computational unit back
to the original station-date identity.

\section{Storage layers and why no single format was sufficient}
\begin{figure}[H]
\centering
\includesvg[inkscapelatex=false,width=\textwidth,height=0.72\textheight,keepaspectratio]{Builtsvg/latex/5}
\caption{The storage stack separates authoritative inputs, scientific arrays, training representations, and evidence products.}
\label{fig:storage-layers}
\end{figure}

Source formats were retained where they carried provenance. SQLite was used for
small mutable state because it supports indexed lookup, transactions, and
backup snapshots. Zarr was used for chunked multidimensional arrays because a
single station-date read should not require loading an entire year. Parquet was
used for tabular joins and audits because columns can be read selectively.
NumPy memory maps were used for training because they permit fixed-offset
access without materialising all rows in RAM. NPZ shards were used for
resumable remote extraction because a completed shard is independently
verifiable and transferable.

The resulting design deliberately contains duplicate representations. That is
not accidental storage waste: each representation is optimised for a different
access pattern. The source layer is authoritative, the array layer is
structured, the cache layer is fast, and the tensor layer is model-shaped.
Hashes, manifests, shapes, and row-order sidecars are what make those
representations equivalent rather than ambiguous copies.

\section{Resumable extraction as a state machine}
\begin{figure}[H]
\centering
\includesvg[inkscapelatex=false,width=\textwidth,height=0.72\textheight,keepaspectratio]{Builtsvg/latex/6}
\caption{A failed remote operation returns to the last trusted state rather than invalidating the complete extraction.}
\label{fig:resume-state}
\end{figure}

The distinction between an attempted operation and a committed operation is
central. A scene search can be repeated safely because its result is cached.
An array write is not considered complete merely because a file exists: the
workflow also records dimensions, valid-value counts, source identifiers, and
a checksum. A partial write therefore cannot masquerade as a valid shard after
a machine interruption. SQLite backup and snapshot operations provide a second
line of defence when the mutable state database itself must be moved between
machines.

\section{Request geometry, compression, and network cost}
Remote raster access has at least four costs: request latency, bytes
transferred, decompression, and local window assembly. A small request is not
always cheaper because authentication and HTTP setup are paid per request. A
large request is not always cheaper because it transfers pixels that will never
enter a sample. The production extraction balanced these costs by clustering
nearby requests, using COG overviews when the target resolution permitted it,
and retaining exact native reads for validation.

The overview path changes the physical pixel support of the input. It is
therefore represented as a named strategy rather than silently replacing the
native product. The output records the strategy, source asset, requested
window, valid-pixel count, and fill behaviour. This turns a performance
approximation into an explicit experimental factor. A future ablation can
compare exact and overview-derived features without reconstructing the entire
archive.

\section{Cache signatures and invalidation}
Persistent caches are useful only if stale values cannot be confused with new
values. Cache names and manifests include the contract version, source
namespace, feature list, scaling policy, precision, clipping policy, and
producer signature. A change to any of these fields produces a new cache
identity. Temporary files are promoted atomically only after validation, so a
process killed during construction leaves the previous complete cache intact.

This design also protects scientific interpretation. If a scaler is changed,
the model input distribution has changed. If a fill value is changed, the
missingness semantics have changed. If a raster row order is changed, the
station-date join may be wrong even though every array has the expected shape.
Signatures make such changes visible in the release manifest instead of
leaving them as undocumented differences between runs.

\section{Memory traffic during training}
Training performance was constrained by memory movement as much as by matrix
multiplication. The cache stores a value channel and a validity channel, uses
FP16 for the normalised local arrays, and gathers only the rows and temporal
steps needed by a batch. Bounded workers and prefetch queues prevent the data
loader from creating an unbounded backlog. Pinned host memory permits
asynchronous transfer to the GPU where the framework and hardware support it.

Gradient accumulation separates the optimisation batch from the physical batch
that fits in device memory. Several small forward/backward passes are summed
before an optimiser update, approximating a larger batch without allocating
all examples simultaneously. Mixed precision reduces activation storage, while
gradient clipping prevents a rare large residual or invalid intermediate from
destabilising an update. These mechanisms address different bottlenecks and
cannot be treated as interchangeable speed settings.

\chapter{Model Architecture and Physical Interpretation}

The model is introduced only after the data contract and computational pathway
are established. This is deliberate: the architecture is a response to the
available measurements and their missingness patterns, not an isolated choice
of neural-network components. Each branch has a defined role, and the fusion
stage combines those roles while retaining enough structure to analyse the
result.

\section{Final v19.5 configuration}
\begin{itemize}
    \item Architecture: multiscale residual spatial-temporal attention gated fusion.
    \item Spatial width: 48.
    \item Temporal hidden width: 64.
    \item Fusion width: 192.
    \item Eight history steps in a maximum 30-day window.
    \item AdamW, MSE in the final exported configuration.
    \item Effective batch size 2,048.
    \item EMA decay 0.85 and gradient clipping norm 2.0.
\end{itemize}

\section{Full model diagram}
\begin{figure}[H]
\centering
\includesvg[inkscapelatex=false,width=\textwidth,height=0.72\textheight,keepaspectratio]{Builtsvg/latex/7}
\caption{v19.5 model architecture. Separate branches preserve measurement semantics before adaptive fusion.}
\end{figure}

\section{Spatial encoders}
EVI, LST, and Sentinel-1 receive independent stems because their values have
different units, noise, missingness, and physical interpretations. Each stem
uses convolution, GroupNorm, GELU, and multi-scale residual blocks. A local
depthwise $3\times3$ path captures nearby structure; a dilated path enlarges
context; a $1\times1$ mixer combines channels; and the residual connection
stabilises optimisation.

Mean pooling describes the average neighbourhood. Max pooling retains a
salient extreme. Attention pooling learns which spatial positions matter for
the current example. Concatenating all three avoids forcing one summary
statistic to represent every landscape.

\section{Temporal TCN}
The temporal branch projects a sequence into hidden width, applies ordinary and
dilated one-dimensional convolutions, scores valid history steps, and gates the
last valid state against the broader history. This resembles a learned
finite-memory filter: it can represent rainfall persistence, delayed
infiltration, evaporation-driven drying, and recent-event dominance.

\section{Static and location branches}
Clay, sand, silt, elevation, slope, and circular aspect are projected through a
static MLP. A location-history configuration also receives a controlled spatial
representation. This can improve regime interpolation but creates a memorisation
risk; a network-disjoint split is stronger than random rows but is not identical
to a continent or biome holdout.

\section{Modality gate and fusion}
The five-way gate produces non-negative weights summing to one over EVI, LST,
Sentinel-1, temporal history, and static context. These weights are predictive
routing coefficients, not causal contribution fractions. The fusion MLP then
models conditional interactions such as precipitation followed by high VPD on
sandy soil or radar structure under dense vegetation.

\section{Calibration}
The calibration layer is:
\[
\hat{y}=c_0+c_1a+c_2r+c_3ar+c_4a^2+c_5r^2,
\]
where $a$ is the anchor and $r$ is the raw model output or residual representation.
Identity initialisation means the layer begins as a sensible no-correction map
and learns only the correction supported by the data.

\feat{A calibration layer can improve output scale without adding new
physical representation. The calibrated score belongs to the complete system,
not necessarily to the uncalibrated neural encoder alone.}

\chapter{Training, Numerical Stability, and Reproducibility}

Training is the point at which the prepared data, model design, and hardware
constraints meet. A model can be mathematically well specified and still
produce an unreliable experiment if batches are too large for memory, if
mixed-precision arithmetic is unstable, or if a checkpoint cannot be traced
back to its input contract. These sections therefore describe the optimisation
procedure together with the safeguards that make it repeatable.

\section{Hyperparameter search}
The project uses staged Optuna studies rather than fully training every
candidate. Shorter and bounded trials screen architectures and optimisation
settings before full retraining. The final exported configuration is treated as
the authority when historical notebook defaults differ.

\section{Loss and optimisation}
The final exported v19.5 configuration records MSE:
\[
\mathcal{L}_{\mathrm{MSE}}=\frac{1}{N}\sum_{i=1}^{N}
(\hat{y}_i-y_i)^2.
\]
AdamW updates the weights using gradients from backpropagation while weight
decay discourages unnecessarily large parameters. Gradient clipping limits
rare extreme updates.

\section{AMP, accumulation, and EMA}
Automatic mixed precision reduces memory and can accelerate matrix operations.
Gradient accumulation approximates a large batch through several microbatches.
EMA maintains a smoothed parameter state that is less sensitive to the noise of
the final individual updates. These mechanisms improve feasibility but add
state and numerical conditions that belong in the reproducibility record.

\section{Compilation and hardware dependence}
The v19.5 training path smoke-tests \texttt{torch.compile} with a
reduce-overhead mode and falls back to eager execution when compilation is
unstable. This is a practical engineering compromise: compilation can reduce
framework overhead but must not be a single point of failure.

\section{Checkpoint policy}
Checkpoints are written atomically, and configuration, scaler, manifest, and
hash information is retained. The historical workflow includes test-$R^2$
checkpoint selection; this is explicitly documented because repeated test
selection means the nominal test set should be interpreted as
development-held-out rather than perfectly untouched.

\chapter{Evaluation and Error Analysis}

The evaluation is presented as a set of connected questions rather than as one
headline number. First, the aggregate scores establish the overall level of
agreement. The diagnostic analyses then ask where the errors occur, whether
they change with time or moisture regime, and whether the apparent gain is
consistent across stations. This order keeps the interpretation close to the
measurement process and avoids treating a single average as a complete
description of model behaviour.

\section{Evaluation philosophy}
Evaluation is treated as an observational campaign analysis. One global score
cannot show whether errors are concentrated in a network, season, moisture bin,
region, acquisition-age stratum, or missingness regime.

\section{Primary results}
\begin{center}
\begin{tabular}{lrrrr}
\toprule
Configuration & Rows & RMSE & MAE & $R^2$\\
\midrule
Environment-only & 644,716 & 0.08626 & 0.06794 & 0.47046\\
Location-history & 905,248 & 0.06594 & 0.04854 & 0.65341\\
\bottomrule
\end{tabular}
\end{center}

The location-history model wins at 62.62\% of matched stations, with paired
RMSE 0.06704 versus 0.08490 for the environment-only model. These paired
statistics are more informative than comparing aggregate row counts alone.

\section{Temporal and regional diagnostics}
Year-by-year performance is strong in 2023--2024 and weaker in 2025--2026.
This may reflect distribution shift, event composition, station mix, coverage,
or target/raster availability. The analysis reports the pattern without
claiming a causal explanation.

\begin{figure}[H]
\centering
\safeimage{error_analysis_results_v19_5/plots/annual_test_rmse.png}{0.78}
\caption{Annual test RMSE diagnostic. The temporal pattern is a monitoring
signal, not by itself a causal explanation.}
\end{figure}

\section{Figures from the global analysis}
\begin{figure}[H]
\centering
\safeimage{error_analysis_results_v19_5/plots/overall_test_mae_rmse.png}{0.72}
\caption{Overall held-out MAE and RMSE comparison.}
\end{figure}

\begin{figure}[H]
\centering
\safeimage{error_analysis_results_v19_5/plots/observed_vs_predicted_density_location_history_test.png}{0.78}
\caption{Observed versus predicted density for the location-history test export.}
\end{figure}

\begin{figure}[H]
\centering
\safeimage{error_analysis_results_v19_5/plots/test_residual_distributions.png}{0.78}
\caption{Residual distributions for the held-out configurations.}
\end{figure}

\section{Oracle trimming}
Oracle trimming ranks predictions by their realised error using the true target
and recomputes retained-set $R^2$. The location-history model reaches $R^2=0.70$
after removing 1.30\% of its worst errors, $R^2=0.80$ after 7.90\%, and
$R^2=0.90$ after 24.10\%. The environment-only model requires 16.80\%, 32.30\%,
and 54.10\%, respectively.

\begin{figure}[H]
\centering
\safeimage{error_analysis_results_v19_5/plots/high_error_trim_r2_sensitivity.png}{0.82}
\caption{Oracle high-error trimming sensitivity. This is a diagnostic of error
concentration, not an operational filtering method.}
\end{figure}

\feat{A deployed system cannot know the true target before prediction.
Operational filtering must use information available before prediction, such as
missingness, acquisition age, uncertainty, or regime flags.}

\chapter{Extended Diagnostic Gallery and Technical Case Studies}

\section{Why the gallery is part of the result}
The figures in this chapter are not decorative appendices. Each plot answers a
different failure question. A global average tests aggregate agreement; a
station-level diagnostic tests concentration across sites; a month-of-year plot
tests seasonal dependence; a moisture-bin plot tests range compression; and a
paired station comparison tests
whether one configuration consistently helps the same locations; residual
plots test whether calibration has removed scale-dependent bias.

\section{Monthly and seasonal behaviour}
\begin{figure}[H]
\centering
\safeimage{error_analysis_results_v19_5/plots/monthly_test_rmse.png}{0.82}
\caption{Monthly test RMSE. Seasonal error is a property of the joint
observation and acquisition system, not automatically a model defect.}
\end{figure}

\begin{figure}[H]
\centering
\safeimage{error_analysis_results_v19_5/plots/month_of_year_test_rmse.png}{0.82}
\caption{Month-of-year test RMSE. This view separates calendar seasonality from
the irregular distribution of individual years.}
\end{figure}

Seasonal variation can arise from changing vegetation, cloud availability,
snow or freeze conditions, rainfall-event frequency, acquisition geometry, and
the depth response of the target. Interpreting a seasonal curve therefore
requires the acquisition-age and missingness diagnostics alongside it.

\section{Moisture range and residual structure}
\begin{figure}[H]
\centering
\safeimage{error_analysis_results_v19_5/plots/test_rmse_by_moisture_bin.png}{0.82}
\caption{RMSE by target-moisture bin. A global score can hide whether the model
performs differently in dry, intermediate, and wet states.}
\end{figure}

\begin{figure}[H]
\centering
\safeimage{error_analysis_results_v19_5/plots/residual_vs_target_location_history_test.png}{0.82}
\caption{Residual versus target for location-history. Curvature or funnel shape
would indicate remaining scale-dependent error.}
\end{figure}

\begin{figure}[H]
\centering
\safeimage{error_analysis_results_v19_5/plots/test_residual_distributions.png}{0.82}
\caption{Location-history residual distribution. The centre, asymmetry, and
tail are distinct pieces of evidence.}
\end{figure}

\section{Raw versus calibrated output}
\begin{figure}[H]
\centering
\safeimage{error_analysis_results_v19_5/plots/raw_vs_calibrated_location_history_test.png}{0.82}
\caption{Raw and calibrated location-history predictions. Calibration is a
measurable part of the production system.}
\end{figure}

\begin{figure}[H]
\centering
\safeimage{error_analysis_results_v19_5/plots/raw_vs_calibrated_density_location_history_test.png}{0.82}
\caption{Density comparison before and after calibration. The visual shift
illustrates why reporting only the final calibrated score is incomplete.}
\end{figure}

The raw-to-calibrated transformation is scientifically informative. A model can
learn useful ranking and conditional structure while producing an output with
incorrect gain or offset. The calibrator can correct this numerical mismatch,
but its coefficients must be fitted on development data and frozen before an
untouched certification test.

\section{Paired station and regional comparisons}
\begin{figure}[H]
\centering
\safeimage{error_analysis_results_v19_5/plots/paired_test_metrics.png}{0.82}
\caption{Paired configuration comparison. Pairing removes some confounding
caused by different station populations.}
\end{figure}

\begin{figure}[H]
\centering
\safeimage{error_analysis_results_v19_5/plots/paired_station_rmse_difference_test.png}{0.82}
\caption{Station-level paired RMSE differences. Improvement is not spatially
uniform, which is why the aggregate result needs station diagnostics.}
\end{figure}

\section{Worst-station and event-context analyses}
\begin{figure}[H]
\centering
\safeimage{error_analysis_results_v19_5/plots/worst_30_stations_location_history_test.png}{0.82}
\caption{Worst 30 location-history stations. These are investigation targets,
not evidence that the whole model fails in the same way.}
\end{figure}

\begin{figure}[H]
\centering
\safeimage{error_analysis_results_v19_5/plots/worst_30_stations_environment_only_test.png}{0.82}
\caption{Worst 30 environment-only stations. Comparing the lists shows whether
location/history changes the identity of difficult regimes.}
\end{figure}

\begin{figure}[H]
\centering
\safeimage{error_analysis_results_v19_5/plots/late_2025_spike_context.png}{0.82}
\caption{Late-2025 event context. Error spikes need event-level inspection
before being attributed to a model architecture.}
\end{figure}

\section{What the gallery establishes}
The diagnostic gallery establishes that the result is not a single accidental
number. It exposes temporal variation, station heterogeneity, range dependence,
calibration dependence, and the long tail of hard cases. It does not by itself
establish causal physical response or area-weighted global accuracy. Those
claims require additional experimental design, especially block uncertainty
and ecological holdouts.

\chapter{Technical Hurdles and Solutions}

\section{Data association was harder than tensor arithmetic}
A tensor can have the correct rank, dtype, and finite values while belonging to
the wrong station and date. The project therefore treated identity as a
scientific product. Station/date keys, cell keys, row indexes, manifests, and
sidecar tables survived across extraction, merging, caching, and training.

\section{Same date is not necessarily same physical state}
An ISMN daily mean and a nearest-valid Sentinel-1 acquisition can be separated
by rainfall, evaporation, irrigation, freezing, or drainage. The resulting
error contains temporal representativeness uncertainty. Strict same-day matching
would reduce this uncertainty but also reduce global coverage.

\section{Calibration can dominate the headline score}
The raw location-history model is poorly scaled relative to the final target,
while calibration makes it useful. The complete system therefore deserves the
calibrated score, but the representation should not be credited with all of the
calibration improvement. Raw and calibrated metrics are both reported.

\section{Location is useful and dangerous}
Coordinates and station history help represent broad climate and soil regimes.
They can also encode site identity. A network-disjoint split limits direct
network leakage, but does not prove transfer to unseen ecological or continental
regimes. This is why the environment-only comparison is a critical control.

\section{Missingness can become a hidden predictor}
Validity masks preserve the difference between a low measurement and no
measurement. However, the mask itself can correlate with season, cloud cover,
region, sensor operation, or network. A model can exploit availability rather
than physical signal. Missingness-stratified performance and mask-only
baselines are appropriate follow-up tests.

\section{Physical validity exceeds numerical validity}
Finite-value and range checks do not prove hydrological plausibility. Strong
aggregate $R^2$ can coexist with implausible wetting transitions, spatial
discontinuities, or climate-specific failures. Event timing, persistence,
moisture-bin residuals, texture-stratified errors, and monotonic stress tests
are needed for stronger physical claims.

\chapter{Reproducibility, Audit Trail, and Release}

\section{What is recorded}
The release surface includes:
\begin{itemize}
    \item contract and feature schema;
    \item train/validation/test manifests;
    \item train-only scaler files and hashes;
    \item cache build signatures;
    \item Zarr and memory-map manifests;
    \item checkpoint configuration and SHA-256 hash;
    \item prediction Parquet outputs;
    \item metric tables and analysis manifest;
    \item exact-versus-overview strategy namespaces;
    \item Azure transfer checksums and completion metadata.
\end{itemize}

These records answer different reproducibility questions. The contract says
what a value is supposed to mean. The split manifest says which identities were
allowed to influence fitting, selection, and testing. The scaler hash says
which numerical transformation was applied before the model saw a value. The
cache signature says which producer, precision, fill convention, and source
namespace created the array. The checkpoint hash says which learned parameters
generated the prediction. Reproducibility is therefore a chain of compatible
identities rather than a single random seed.

The chain is also designed for diagnosis. If a metric changes, the audit trail
allows the change to be narrowed to the source catalog, daily aggregation,
remote asset selection, raster alignment, cache construction, model state, or
metric script. Without this separation, a changed score would be visible but
its cause would not be. That distinction is especially important in a
multi-provider experiment where a source catalog or remote asset can change
independently of the model code.

\section{Failure should be explicit}
The builders reject missing columns, duplicate station/date rows, unsupported
split labels, missing features, wrong shapes, mixed temporal lengths, non-finite
arrays, and invalid numeric targets. This is preferable to silently coercing
bad input into a success-shaped training set.

The checks are placed at boundaries because errors become harder to localise
after data have been merged. At the archive boundary, malformed files and
unsupported depth metadata are rejected. At the daily-target boundary, duplicate
identities and invalid quality flags are exposed. At the raster boundary,
shape, row-order, source-asset, and valid-pixel checks are applied. At the
tensor boundary, feature dimensions, masks, split IDs, and finite values are
checked. At the evaluation boundary, prediction counts and key uniqueness are
compared with the expected manifest.

This boundary design also distinguishes absence from failure. A missing remote
observation can be a valid scientific condition and is represented by a mask.
A missing required shard, a mismatched row count, or a stale cache signature is
a pipeline failure and must stop the run. Treating both cases as a generic zero
would make the model appear complete while changing the physical meaning of
the input.

\section{Reproduction sequence}
\begin{enumerate}
    \item Verify the contract and source manifests.
    \item Rebuild or validate the catalog and daily Parquet.
    \item Verify ERA5, SMAP, MODIS, Sentinel-1, and static-product provenance.
    \item Validate cluster inventory, shards, checksums, and Zarr metadata.
    \item Reconstruct train-only scalers and cache signatures.
    \item Validate sample IDs, row order, shapes, and finite values.
    \item Restore the exact checkpoint and model configuration.
    \item Recompute predictions and compare full-count metric tables.
\end{enumerate}

The sequence is intentionally ordered from meaning to computation. Rebuilding
the neural network before confirming the target contract would reproduce a
number without proving that it refers to the same task. Rebuilding the cache
before checking the producer row order could reproduce a tensor shape while
silently permuting the station associations. The first stages are therefore
slower and more administrative than model loading, but they prevent a fast
reproduction of the wrong experiment.

For a complete rerun, each stage should emit a compact manifest containing
inputs, outputs, counts, hashes, elapsed time, and the reason for any skipped
work. A skipped unit is reproducible only when its prior completion record is
still valid under the current contract and source namespace. Otherwise it must
be rebuilt. This is the operational meaning of resumability: reuse is allowed
only when equivalence has been demonstrated.

\section{Change control across providers}
The data providers evolve on different schedules. A remote catalogue can expose
a new asset while the station archive remains unchanged; a processing library
can alter how a COG window is decoded; and a model notebook can change its
default precision without changing its filename. The release process therefore
records provider, collection, asset identifier, acquisition date, access
strategy, and producer version separately. A filename alone is not considered
an adequate provenance key.

When a source changes, the correct response depends on the layer affected. A
new static soil layer requires rebuilding the static feature table but not
necessarily the dynamic scene inventory. A new Sentinel asset can require
re-running only the clusters whose search results changed. A scaler or fill
policy change invalidates the normalised cache and all dependent tensors, even
if the source arrays are byte-for-byte identical. This dependency structure
prevents both unnecessary global rebuilds and unsafe partial reuse.

The same principle applies to software. A numerical library update is tested
first on a bounded sample with known row counts, valid-pixel counts, shapes,
and summary statistics. Only after those invariants match is the new producer
allowed to write a production namespace. This is the data-engineering
equivalent of checking a new numerical integrator against a solved test
problem before using it in a long simulation.

\section{Why auditability matters scientifically}
Auditability is not only a software quality property. It determines whether a
physics interpretation can be challenged and reproduced. If a location-history
model improves over an environment-only model, the reader needs to know
whether the improvement came from temporal memory, local static context,
station identity, a calibration transformation, or a difference in available
rows. The manifests and paired evaluation tables make those alternatives
separable.

The same applies to a failure. A low score may reflect poor representation,
temporal mismatch, missingness, a remote extraction gap, or a target aggregation
problem. A reproducible pipeline does not guarantee that the model is correct,
but it ensures that competing explanations can be tested against the same
intermediate artifacts. That is the practical meaning of making a global
machine-learning experiment scientifically inspectable.

\chapter{Conclusion}

The global soil-moisture project is a substantial technical accomplishment because it
integrates scientific data interpretation, cloud-scale acquisition, storage
engineering, numerical optimisation, machine learning, and evaluation into one
auditable system. The model score is only the visible endpoint. The deeper
result is that a global experiment involving heterogeneous station
observations, multi-resolution remote products, reanalysis, static layers,
irregular histories, and finite compute was made operationally possible.

The engineering strategy is coherent:
\begin{enumerate}
    \item define the physical target before extracting predictors;
    \item preserve identity and provenance through every representation;
    \item share remote work through clustering and reuse;
    \item make interrupted work resumable and checksummed;
    \item move deterministic preparation into persistent caches;
    \item train a modality-aware model under bounded memory;
    \item evaluate at full population scale while exposing limitations.
\end{enumerate}

The final location-history configuration achieves calibrated RMSE
$0.06594$ and $R^2=0.65341$ over 905,248 test rows, compared with RMSE
$0.08626$ and $R^2=0.47046$ for the environment-only configuration. These
results are meaningful because they are attached to a documented data contract,
explicit provenance, extensive error analysis, and a clear statement of what
the scores do and do not establish.

\feat{When the unoptimised experiment is impossible, optimisation becomes part
of the scientific methodology. Here, reducing redundancy without erasing
provenance converted computational feasibility into repeatable evidence.}

\chapter{Future Work}

\section{Certification of physical transfer}
The strongest next step is a set of evaluation protocols that separate
interpolation among familiar regimes from transfer to genuinely novel regimes.
The existing network-disjoint split is valuable, but future experiments should
also hold out continents, climate zones, biomes, soil-texture regimes, and
station density classes. These tests would show whether the location-history
branch is learning transferable environmental structure or primarily regional
priors.

\section{Depth-specific and temporally synchronous targets}
Headline metrics should be separated for 0--5~cm and 5--10~cm observations.
The predictor-target acquisition age should be recorded as a first-class
feature, and errors should be stratified by the offset between the target day
and Sentinel-1 or MODIS support window. A strict same-day sensitivity experiment
would quantify the price of temporal synchrony in coverage.

\section{Uncertainty and statistically dependent observations}
Station time series are autocorrelated, so confidence intervals should be
estimated by station or network blocks rather than by treating every row as
independent. Conformal prediction, ensemble dispersion, or a calibrated
heteroscedastic head could provide uncertainty intervals. These intervals
should be evaluated for coverage during missing-data, event, and acquisition-age
regimes.

\section{Physical stress tests}
Future analysis should test wetting after precipitation, drying under high VPD,
persistence across short gaps, response across soil texture, and behaviour near
dry and wet bounds. These tests would not impose an unjustified deterministic
law on noisy observations; they would identify clearly implausible response
patterns that aggregate regression scores can conceal.

\section{Representation and approximation ablations}
The exact native Sentinel-1 path should be compared directly with the
overview-derived path. FP32 and unclipped cache variants should be evaluated
against the production FP16 representation. Modality ablations should remove
EVI, LST, radar, temporal history, static context, location, masks, and the
SMAP anchor one at a time. This would quantify which technical approximations
and information pathways are responsible for the final improvement.

\section{Locked certification release}
The final certification workflow should fit all scalers, calibration
coefficients, hyperparameters, and checkpoint choices using training and
validation only, then evaluate once on a locked test cohort. The current report
records the historical test-$R^2$ selection policy transparently; a future
locked release would provide a cleaner estimate of untouched generalisation.

\section{Operational data system}
The resumable acquisition system can be extended with automatic provenance
lineage, remote asset version snapshots, formal schema migrations, and
hardware-specific throughput dashboards. The goal would be to make every
production prediction traceable to a source asset, acquisition time, cache
signature, model checkpoint, and calibration version.

\clearpage
\renewcommand{\bibname}{References}
\nocite{*}
\bibliographystyle{plainnat}
\bibliography{references}

\end{document}
